\documentclass[conference]{IEEEtran}

\usepackage{iftex}
\ifPDFTeX
\usepackage[utf8]{inputenc}
\usepackage[T1]{fontenc}
\fi

\usepackage{amsmath,bm,mathtools}
\usepackage{newtxtext,newtxmath}
\usepackage{textcomp}
\usepackage{enumitem}
\usepackage{booktabs}
\usepackage{siunitx}
\usepackage{float}
\usepackage{pgf}
\usepackage{pgfplots}
\usepackage{tikz}
\usepackage[breaklinks=true,hidelinks]{hyperref}
\usepackage[nameinlink,capitalise,noabbrev]{cleveref}
\usepackage[final]{microtype}
% Matplotlib PGF backends may emit \mathdefault; define fallback for LuaLaTeX.
\providecommand{\mathdefault}[1]{#1}

% For including PGF
\newcommand{\IncludePGF}[1]{\resizebox{\linewidth}{!}{\input{#1}}}

\interdisplaylinepenalty=2500
\allowdisplaybreaks

\DeclareMathOperator{\clamp}{clamp}
\DeclareMathOperator{\atanTwo}{atan2}

\newcommand{\vct}[1]{\bm{#1}}
\newcommand{\R}{\mathbb{R}}
\newcommand{\dd}{\,\mathrm{d}}

\title{Adaptive Vertical PII with Mahony Attitude and Frequency-Based Scheduling\\
(Implementation Note)}

\author{%
    \IEEEauthorblockN{Mikhail S.\ Grushinskiy}
    \IEEEauthorblockA{Bareboat Necessities Project\\
    \texttt{mgrouch@users.github.com}}
}

\begin{document}
    \maketitle

    \begin{abstract}
        This note documents the method implemented by
        \texttt{AdaptiveVerticalPIIMahony.h},
        \texttt{AdaptiveVerticalPII.h}, and
        \texttt{VerticalPIIObserver.h}.
        The pipeline estimates attitude with a Mahony complementary filter,
        rotates body-frame specific force into a \(Z\)-up world frame,
        subtracts gravity to obtain vertical inertial acceleration,
        and drives an adaptive vertical displacement observer whose gains are
        scheduled from a frequency proxy and a sea-state amplitude estimate.

        The implementation is layered.
        First, a repeated-real-pole vertical PII observer estimates
        displacement and velocity from world-frame vertical inertial acceleration.
        Second, an adaptive wrapper adds an acceleration-sigma estimator, an
        OU-style sea-state envelope estimator, and a policy-selected frequency
        tracker, and schedules the observer parameters from them.
        Third, a Mahony wrapper supplies the vertical acceleration from raw IMU
        data and schedules its own attitude gains from accelerometer trust.
        The emphasis here is on the actual implemented equations and signal flow.
    \end{abstract}

    \begin{IEEEkeywords}
        heave estimation, vertical observer, PII observer, Mahony AHRS,
        adaptive scheduling, marine inertial sensing
    \end{IEEEkeywords}

% ---------------------------------------------------------------
    \section{Signal Flow and Layering}

    The implementation is organized as a small stack:
    \begin{enumerate}[leftmargin=1.5em]
        \item \texttt{VerticalPIIObserver}: core vertical-motion observer,
        \item \texttt{AdaptiveVerticalPII}: wrapper adding sigma estimation, an
        OU-style sea-state envelope estimator, frequency-proxy scheduling, and
        tracker-policy abstraction,
        \item \texttt{AdaptiveVerticalPIIMahony}: wrapper that uses Mahony AHRS
        to convert raw IMU data into vertical inertial acceleration.
    \end{enumerate}

    The simulation entry point instantiates
    \texttt{AdaptiveVerticalPIIMahony<float,\ true,\ TrackerType::PLL>}
    in the supplied driver, but the adaptive wrapper itself is generic and can
    use several tracker backends through \texttt{FrequencyTrackerPolicy}.

    At sample time \(k\) with period \(\Delta t_k\), the implemented cascade is
    \begin{align}
        \vct{\omega}_{b,k} &\xrightarrow{\text{Mahony}} \hat{q}_k, \\
        \vct{f}_{b,k} &\xrightarrow{\hat{q}_k} \hat{\vct{f}}_{w,k}, \\
        \hat{a}_{z,k} &= \hat{f}_{w,z,k} - g, \\
        \hat{a}_{z,k} &\xrightarrow{\text{Adaptive Vertical PII}}
        (\hat{p}_k,\hat{v}_k,\hat{S}_k,\hat{b}_k).
    \end{align}
    Here:
    \begin{itemize}[leftmargin=1.4em]
        \item \(\hat{q}\) is the stored Mahony quaternion, exposed as
        \texttt{quaternionWorldToBody()},
        \item \(\hat{p}\) is vertical displacement,
        \item \(\hat{v}\) is vertical velocity,
        \item \(\hat{S}\) is the integral of displacement,
        \item \(\hat{b}\) is the optional slow acceleration-bias estimate.
    \end{itemize}

    Inside \texttt{AdaptiveVerticalPII::update()} the per-sample order is fixed:
    \begin{enumerate}[leftmargin=1.5em]
        \item update the running acceleration mean / variance / sigma,
        \item step the frequency tracker with \(\hat{a}_z/g\),
        \item update the sea-state envelope estimator,
        \item accumulate \(dt\) and, once
        \(\sum dt \ge\) \texttt{auto\_schedule\_period\_s}, run one scheduling
        update,
        \item run the core observer for this sample.
    \end{enumerate}

    The rest of this note follows the code structure:
    first the vertical PII observer, then the adaptive scheduling layer and its
    envelope estimator, then the frequency proxy, then the Mahony wrapper.

% ---------------------------------------------------------------
    \section{Vertical PII Observer}\label{sec:pii}

    \subsection{Purpose and States}

    \texttt{VerticalPIIObserver} estimates vertical displacement and velocity from
    an input that is already the gravity-compensated world-frame vertical inertial
    acceleration.

    Its main states are:
    \begin{align}
        a_f &: \text{filtered vertical acceleration},\\
        v   &: \text{vertical velocity},\\
        p   &: \text{vertical displacement},\\
        S   &: \text{integral of displacement}.
    \end{align}

    When the optional bias channel is enabled, it also keeps
    \begin{align}
        d &: \text{very slow displacement trend},\\
        b &: \text{estimated acceleration bias}.
    \end{align}

    The bias channel is a compile-time option (\texttt{WithBias}); with
    \texttt{WithBias=false} the two extra states are not even instantiated and
    the accessors return zero.

    \subsection{Continuous-Time Model}

    The observer equations documented in the header are
    \begin{align}
        \dot{a}_f &= \frac{a_{\mathrm{meas}} - a_f}{\tau_a}, \label{eq:afdot}\\
        \dot{S}   &= p, \label{eq:Sdot}\\
        \dot{p}   &= v, \label{eq:pdot}\\
        \dot{v}   &= a_f - b - k_v v - k_p p - k_s S. \label{eq:vdot}
    \end{align}

    When \texttt{WithBias=true}, the very-slow bias-trend channel is
    \begin{align}
        \dot{d} &= \frac{p-d}{\tau_d}, \label{eq:ddot}\\
        \dot{b} &= k_b d - \lambda_b b. \label{eq:bdot}
    \end{align}

    So the observer is not just a double integrator. It is a damped
    third-order restoring system driven by filtered acceleration, with an
    optional very-slow leak-stabilized bias estimator.

    \subsection{Repeated-Real-Pole Parameterization}

    The core PII restoring loop is parameterized by a repeated real pole at
    \(-r\):
    \begin{equation}
    (s+r)^3 = s^3 + 3rs^2 + 3r^2 s + r^3.
    \end{equation}

    Matching coefficients yields
    \begin{equation}
        k_v = 3r,\qquad
        k_p = 3r^2,\qquad
        k_s = r^3.
    \end{equation}

    This matters because the observer can be tuned through one physically
    meaningful bandwidth-like parameter \(r\), while preserving an overdamped,
    all-real-pole structure. The implementation derives the active gains
    \((k_v,k_p,k_s)\) from the currently active \(r\) every time the active
    motion-pole rate changes.

    There is also a manual escape hatch, \texttt{set\_active\_pii\_gains()},
    which writes \((k_v,k_p,k_s)\) directly. Because that breaks the
    repeated-pole relation, it disables adaptation as a side effect.

    \subsection{Discrete Update Used by the Code}

    The implementation uses the embedded-friendly one-pole coefficient
    \begin{equation}
        \alpha(dt,\tau) = \frac{dt}{\tau + dt},
    \end{equation}
    rather than an exponential form.

    The filtered acceleration update is
    \begin{equation}
        a_f^{+} = a_f + \alpha(dt,\tau_a)(a_{\mathrm{meas}}-a_f).
    \end{equation}
    A non-finite \(a_{\mathrm{meas}}\) is replaced by the current \(a_f\), which
    makes the sample a hold rather than a fault.

    If the bias channel is enabled, the slow trend update is
    \begin{equation}
        d^{+} = d + \alpha(dt,\tau_d)(p-d),
    \end{equation}
    and the bias state is integrated as
    \begin{equation}
        b^{+} = b + dt\,(k_b d^{+} - \lambda_b b).
    \end{equation}

    The vertical PII core then forms
    \begin{equation}
        \dot{v} = a_f^{+} - b^{+} - k_v v - k_p p - k_s S,
    \end{equation}
    and applies a semi-implicit update:
    \begin{align}
        v^{+} &= v + dt\,\dot{v}, \\
        p^{+} &= p + dt\,v^{+}, \\
        S^{+} &= S + dt\,p^{+}.
    \end{align}

    That update order is important: the new velocity is used immediately in the
    position update, and the new position is used immediately in the integral
    update.

    \subsection{State and Safety Clamps}

    The implementation supports optional symmetric clamps on
    \begin{equation}
        a_f,\quad v,\quad p,\quad S,\quad d,\quad b.
    \end{equation}
    A non-positive limit disables the corresponding clamp. These are practical
    safety bounds rather than part of the nominal observer model.

    \begin{table}[t]
        \centering
        \caption{Main states of \texttt{VerticalPIIObserver}}
        \label{tab:pii_states}
        \small
        \begin{tabular}{@{}lll@{}}
            \toprule
            Quantity & Code field & Meaning \\
            \midrule
            \(a_f\) & \texttt{a\_f\_} & filtered vertical acceleration \\
            \(v\)   & \texttt{v\_}    & vertical velocity \\
            \(p\)   & \texttt{p\_}    & vertical displacement \\
            \(S\)   & \texttt{S\_}    & integral of displacement \\
            \(d\)   & \texttt{bias\_.d} & slow displacement trend \\
            \(b\)   & \texttt{bias\_.b} & estimated acceleration bias \\
            \(r\)   & \texttt{active\_r\_} & active repeated pole rate \\
            \bottomrule
        \end{tabular}
    \end{table}

% ---------------------------------------------------------------
    \section{Adaptive Scheduling Wrapper}\label{sec:adaptive}

    \subsection{Role of \texttt{AdaptiveVerticalPII}}

    \texttt{AdaptiveVerticalPII} wraps the core observer and adds:
    \begin{itemize}[leftmargin=1.4em]
        \item a running estimate of vertical-acceleration mean, variance, and sigma,
        \item an OU-style sea-state envelope estimator producing a correlation
        time \(\tau_{\mathrm{env}}\) and an amplitude \(\sigma_{\mathrm{env}}\),
        \item a policy-selected acceleration-frequency tracker,
        \item automatic or external scheduling calls into
        \texttt{VerticalPIIObserver::update\_adaptation()}.
    \end{itemize}

    It therefore sits between the raw vertical acceleration signal and the
    observer's scheduling interface.

    \subsection{Acceleration Variance Estimate}

    At IMU rate, the wrapper updates a slow mean and a shorter-time variance
    using the same one-pole coefficient \(\alpha(dt,\tau)\):
    \begin{align}
        \mu_a^{+} &= \mu_a + \alpha(dt,\tau_\mu)(a_z-\mu_a),\\
        e_a &= a_z - \mu_a^{+},\\
        \nu_a^{+} &= \nu_a + \alpha(dt,\tau_\nu)(e_a^2 - \nu_a),\\
        \sigma_a^{+} &= \sqrt{\max(\nu_a^{+},0)},
    \end{align}
    with \(\tau_\mu = \SI{20}{s}\) and \(\tau_\nu = \SI{6}{s}\) by default.

    A configured floor is then enforced:
    \begin{equation}
        \sigma_a^{+} \leftarrow \max(\sigma_a^{+}, \sigma_{\min}),
        \qquad \sigma_{\min} = 10^{-4}.
    \end{equation}

    This \(\sigma_a\) is the fallback roughness indicator for observer
    scheduling, and it is also the one used by the Mahony trust adaptation of
    \cref{sec:mahony}.

    \subsection{OU-Style Sea-State Envelope}\label{sec:envelope}

    The wrapper additionally runs an envelope estimator aligned with the
    \texttt{SeaStateAutoTuner} / OU parameterization used elsewhere in this code
    base. It is the preferred source of the scheduling amplitude, and its
    outputs are what the plots label as \(\tau\) and \(\sigma\).

    All four internal averages are debiased exponential moving averages: each
    keeps a value and an accumulated weight and reports value/weight, so the
    estimate is unbiased from the first sample instead of ramping up from zero.
    Unlike the rest of the stack, the envelope uses the exponential coefficient
    \begin{equation}
        \alpha_{\exp}(dt,\tau) = 1 - e^{-dt/\tau}.
    \end{equation}

    The frequency input is the scheduled tracker frequency of
    \cref{sec:freq_proxy}, clamped to \([0.05,5.0]~\si{Hz}\) and smoothed with
    \(\tau_f = \SI{1}{s}\):
    \begin{equation}
        f_{\mathrm{env}} = \mathrm{EMA}\bigl(\clamp(f_{\mathrm{sched}},0.05,5.0),
        \alpha_{\exp}(dt,\tau_f)\bigr).
    \end{equation}

    The variance horizon is then rebuilt from that frequency every sample,
    \begin{equation}
        \tau_{\mathrm{var}} =
        \clamp\!\left(\frac{K_{\mathrm{per}}}{f_{\mathrm{env}}},\,0.30,\,60\right)~\si{s},
        \qquad K_{\mathrm{per}} = 2,
    \end{equation}
    so the averaging window is a fixed number of wave periods rather than a
    fixed number of seconds. With \(\alpha_v=\alpha_{\exp}(dt,\tau_{\mathrm{var}})\),
    \begin{align}
        \mu &= \mathrm{EMA}(a_z,\alpha_v),\\
        \overline{a^2} &= \mathrm{EMA}(a_z^2,\alpha_v),\\
        \nu_{\mathrm{inst}} &= \max\bigl(\overline{a^2}-\mu^2,\,0\bigr),\\
        \nu_{\mathrm{env}} &= \mathrm{EMA}(\nu_{\mathrm{inst}},\alpha_v).
    \end{align}

    The sensor noise floor is removed in variance, and the remainder is mapped
    to the OU amplitude and correlation time:
    \begin{align}
        \sigma_{\mathrm{env}} &=
        \clamp\!\Bigl(c_\sigma\sqrt{\max(\nu_{\mathrm{env}}-\sigma_n^2,0)},\,
        0,\,\sigma_{\max}\Bigr),\\
        \tau_{\mathrm{env}} &=
        \clamp\!\left(\frac{c_\tau}{2 f_{\mathrm{env}}},\,
        \tau_{\min},\,\tau_{\max}\right),
    \end{align}
    with \(\sigma_n=\SI{0.12}{m/s^2}\), \(c_\sigma=0.90\), \(c_\tau=1.38\),
    \(\sigma_{\max}=6\), and \(\tau\in[0.02,3.00]~\si{s}\).
    Both targets are then smoothed toward the applied values with
    \(\alpha_{\exp}(dt,\SI{1.8}{s})\), which is what
    \texttt{envelopeTauApplied()} and \texttt{envelopeSigmaApplied()} return.

    The \(\tau_{\max}\) clamp is active over most of the wave band: since
    \(\tau_{\mathrm{env}}=0.69/f_{\mathrm{env}}\), it saturates at \SI{3}{s} for
    every \(f_{\mathrm{env}}\) below about \SI{0.23}{Hz}, so in typical seas
    \(\tau_{\mathrm{env}}\) is a constant and \(\sigma_{\mathrm{env}}\) carries
    the sea-state information.

    From the same pair the wrapper reports two derived scales,
    \begin{align}
        \text{displacement scale} &= \tfrac12 C_{H_s}\,
        \sigma_{\mathrm{env}}\,\tau_{\mathrm{env}}^2,
        \qquad C_{H_s} = \frac{2\sqrt{2}}{\pi^2},\\
        \text{vertical speed envelope} &= \frac{\sqrt{2}}{\pi}\,
        \sigma_{\mathrm{env}}\,\tau_{\mathrm{env}},
    \end{align}
    both of which return NaN rather than a number when the envelope inputs are
    not valid. The envelope is reported ready only once the frequency and
    variance averages carry weight and all outputs are finite; until then the
    scheduler falls back to \(\sigma_a\).

    \subsection{Three Ways to Drive Adaptation}

    The wrapper exposes three entry points, which differ only in where the
    frequency and sigma come from:
    \begin{itemize}[leftmargin=1.4em]
        \item \texttt{updateAdaptationFromDisplacementFrequency(f,dt,c)}:
        the caller supplies the displacement frequency; sigma is taken from the
        envelope when ready and from \(\sigma_a\) otherwise.
        \item \texttt{updateAdaptationExternal(f,sigma,dt,c)}:
        the caller supplies both. A finite external sigma inside
        \([\sigma_{a,\min},\sigma_{a,\max}]\) wins; otherwise the same
        envelope/\(\sigma_a\) fallback applies.
        \item \texttt{updateAdaptationFromAccelFrequencyProxy(dt)}:
        fully internal. Frequency is the envelope frequency when the envelope is
        ready and the scheduled tracker frequency otherwise; sigma follows the
        same preference.
    \end{itemize}
    A non-finite confidence argument on the two external paths is replaced by
    \texttt{default\_external\_confidence} (\(1.0\)). The internal path computes
    its own confidence as described in \cref{sec:freq_proxy}. Frequencies and
    sigmas taken from the internal sources are clamped into the observer's
    configured validity ranges before being passed down.

    When \texttt{auto\_schedule\_from\_accel\_freq} is enabled, the third entry
    point is called from \texttt{update()} once per
    \texttt{auto\_schedule\_period\_s} (\SI{0.50}{s}), using the accumulated
    time as the scheduling \(dt\).

    \subsection{Scheduling Hook Inside \texttt{VerticalPIIObserver}}

    Whenever adaptation is enabled and a scheduling update is triggered,
    \texttt{VerticalPIIObserver::update\_adaptation()} receives:
    \begin{equation}
    (f_{\mathrm{disp}},\sigma_a,c,dt_{\mathrm{sched}}).
    \end{equation}

    The active parameters are held unchanged, and nothing is smoothed, if
    \(c < c_{\min} = 0.22\), or if \(f_{\mathrm{disp}}\) falls outside
    \([0.02,1.50]~\si{Hz}\), or if \(\sigma_a\) falls outside
    \([10^{-4},50]\), or if \(dt_{\mathrm{sched}}\) is not positive and finite.

    Otherwise the observer first smooths the tracker inputs:
    \begin{align}
        f_f^{+} &= f_f + \alpha_{\mathrm{in}}(f_{\mathrm{disp}}-f_f),\\
        \sigma_f^{+} &= \sigma_f + \alpha_{\mathrm{in}}(\sigma_a-\sigma_f).
    \end{align}

    Then it forms normalized scheduling ratios
    \begin{equation}
        \rho_f = \clamp\!\left(\frac{f_f}{f_{\mathrm{ref}}},0.25,4\right),
        \qquad
        \rho_\sigma = \clamp\!\left(\frac{\sigma_{\mathrm{ref}}}{\sigma_f},0.25,4\right).
    \end{equation}
    Note the inversion: \(\rho_\sigma\) is large in calm water and small in rough
    water.

    The commanded observer parameters are conservative power-law schedules:
    \begin{align}
        r_{\mathrm{cmd}} &=
        r_0\,\rho_f^{\eta_r}\,\rho_\sigma^{\zeta_r},\\
        \tau_{a,\mathrm{cmd}} &=
        \tau_{a,0}\,\rho_f^{\eta_a}\,\rho_\sigma^{\zeta_a},\\
        \tau_{d,\mathrm{cmd}} &=
        \tau_{d,0}\,\rho_f^{\eta_d}\,\rho_\sigma^{\zeta_d},\\
        k_{b,\mathrm{cmd}} &=
        k_{b,0}\,\rho_f^{\eta_b}\,\rho_\sigma^{\zeta_b},
    \end{align}
    where the base values \((r_0,\tau_{a,0},\tau_{d,0},k_{b,0})\) are the
    configured reference parameters, not the currently active ones, so the
    schedule cannot ratchet. Each commanded parameter is clamped to configured
    bounds, then the active parameters move toward the commands via one-pole
    smoothing:
    \begin{equation}
        \theta^{+}
        =
        \theta + \alpha_p(\theta_{\mathrm{cmd}}-\theta),
    \end{equation}
    for each scheduled parameter \(\theta\), after which \(k_v,k_p,k_s\) are
    rebuilt from the new active \(r\). The two bias-channel schedules exist only
    when \texttt{WithBias=true}.

    \subsection{Interpretation of the Scheduling Signs}

    With the exponents of \cref{tab:exponents}, the intended control effect is:
    \begin{itemize}[leftmargin=1.4em]
        \item higher dominant displacement frequency allows a somewhat faster
        observer core (\(\eta_r>0\)) and less acceleration smoothing
        (\(\eta_a<0\)),
        \item larger acceleration sigma reduces restoring authority
        (\(\zeta_r>0\) acting on the inverted ratio) and increases smoothing
        (\(\zeta_a<0\)),
        \item larger sigma also lengthens \(\tau_d\) slightly and shrinks
        \(k_b\), i.e.\ rough water slows bias learning.
    \end{itemize}

    Sigma carries most of the scheduling weight, and not because frequency
    matters less physically. The two are competing estimates of the same sea
    state, and the accel-derived frequency proxy is compressed: the acceleration
    spectrum is weighted by \(\omega^4\), so its dominant frequency sits well
    above the wave peak and barely separates a \SI{3}{s} sea from an \SI{11}{s}
    one. Accel sigma separates the same sea states by a factor of five, so it is
    the better-conditioned regressor of the two.

    The input and parameter smoothing time constants are on the order of a
    minute rather than a few seconds, deliberately. What the schedule tracks is
    the sea state, which moves over tens of minutes; anything faster only lets
    per-wave jitter in the frequency and sigma estimates modulate the observer
    gains, and that modulation shows up directly as heave error.

% ---------------------------------------------------------------
    \section{Frequency Proxy Used for Scheduling}\label{sec:freq_proxy}

    \subsection{Generic Policy Interface}

    \texttt{AdaptiveVerticalPII} does not hard-code one tracker algorithm.
    Instead it uses a policy-selected backend
    \begin{equation}
        \texttt{AccelFreqTracker = TrackerPolicy<TT>},
    \end{equation}
    where \texttt{TT} is a template parameter such as PLL, Aranovskiy, KalmANF,
    or ZeroCross. Convenience aliases
    \texttt{AdaptiveVerticalPII\_PLL} and friends bind each choice.

    Every tracker is fed the same signal the observer sees, scaled to
    dimensionless \(g\) units:
    \begin{equation}
        u_k = \hat{a}_{z,k}/g.
    \end{equation}

    Through the policy helpers, the wrapper queries:
    \begin{itemize}[leftmargin=1.4em]
        \item main frequency estimate,
        \item raw frequency estimate,
        \item tracker confidence,
        \item lock status,
        \item optional coarse estimate.
    \end{itemize}
    The helpers are all \texttt{if constexpr} shims, so a tracker that has no
    \texttt{Config}, no \texttt{reset(f)}, or no coarse channel still compiles.

    \subsection{Scheduled Frequency}

    The scheduled frequency is built by:
    \begin{enumerate}[leftmargin=1.5em]
        \item preferring the tracker's main frequency estimate,
        \item falling back to its raw frequency estimate if that is not positive
        and finite,
        \item blending toward a coarse estimate when the tracker exposes one,
        \begin{equation}
            f_{\mathrm{sched}} \leftarrow
            f_{\mathrm{sched}} + \beta\,(f_{\mathrm{coarse}}-f_{\mathrm{sched}}),
        \end{equation}
        with \(\beta = 0.48\) nominally,
        \item raising that blend to \(\beta \ge 0.75\) whenever the tracker is
        not locked, so an unlocked tracker is mostly overridden by the coarse
        estimate.
    \end{enumerate}

    So the scheduling frequency is intentionally more conservative than just the
    raw tracker output. The same quantity is the frequency input of the envelope
    estimator of \cref{sec:envelope}.

    \subsection{Confidence Passed to the Scheduler}

    The confidence passed to the observer scheduler begins with the tracker's own
    confidence, then is floored upward:
    \begin{align}
        c &\leftarrow \max\bigl(c,\;\max(0.52,\,0.62)\bigr)
        && \text{if a coarse estimate exists},\\
        c &\leftarrow \max(c,\,0.82)
        && \text{if the tracker reports locked},
    \end{align}
    and finally clamped to \([0,1]\).

    This lets the observer continue adapting moderately even when the internal
    frequency proxy is only coarse-valid rather than fully locked. Both floors
    sit above the \(c_{\min}=0.22\) gate, so a coarse-valid proxy is always
    enough to keep the schedule moving.

% ---------------------------------------------------------------
    \section{Mahony Attitude Wrapper}\label{sec:mahony}

    \subsection{World-Frame Convention}

    The supplied Mahony implementation is not NED. In this code base it is used
    as a \(Z\)-up auxiliary/world frame. At identity quaternion, the expected
    accelerometer reading is approximately \((0,0,+g)\), so world up is \(+Z\).

    Accordingly, the wrapper computes
    \begin{equation}
        a_{z,\mathrm{world,up}}
        =
        \bigl(R(\hat{q})\, a_b\bigr)_z - g,
    \end{equation}
    where:
    \begin{itemize}[leftmargin=1.4em]
        \item \(a_b\) is body-frame specific force in \si{m/s^2},
        \item \(R(\hat{q})\) is the rotation built from the stored Mahony
        quaternion,
        \item \(g\) is the configured gravity magnitude.
    \end{itemize}

    This vertical inertial acceleration is exactly the signal fed into the
    adaptive vertical observer; a non-finite result is replaced by zero.

    \subsection{Mahony Innovation Structure}

    The Mahony state stores the quaternion
    \begin{equation}
        \hat{q} = [q_0\;q_1\;q_2\;q_3]^\top .
    \end{equation}

    In the IMU-only branch, the accelerometer is normalized and the predicted
    body-frame direction of world \(+Z\) is
    \begin{equation}
        \hat{\vct{g}}_b =
        \begin{bmatrix}
            2(q_1 q_3 - q_0 q_2)\\
            2(q_0 q_1 + q_2 q_3)\\
            q_0^2 - q_1^2 - q_2^2 + q_3^2
        \end{bmatrix}.
    \end{equation}

    The Mahony routines carry this as a half-vector,
    \(\vct{v}_{1/2}=\hat{\vct{g}}_b/2\), and form the half-innovation
    \begin{equation}
        \vct{e}_{1/2} = \hat{\vct{a}}_b \times \vct{v}_{1/2},
    \end{equation}
    where \(\hat{\vct{a}}_b\) is the normalized measured accelerometer
    direction. The same \(\hat{\vct{g}}_b\) expression, in full-vector form, is
    what the wrapper uses for its trust metric below.

    With magnetometer aiding, the reference field \(\vct{b}=[b_x\;0\;b_z]^\top\)
    is recomputed each sample from the measured field, its predicted body
    direction \(\vct{w}_{1/2}\) is formed the same way, and the innovation
    becomes
    \begin{equation}
        \vct{e}_{1/2}
        = \hat{\vct{a}}_b \times \vct{v}_{1/2}
        + \hat{\vct{m}}_b \times \vct{w}_{1/2}.
    \end{equation}
    An all-zero magnetometer sample falls back to the IMU-only branch.

    \subsection{Mahony Kinematics}

    The corrected rate and its integral term are
    \begin{align}
        \vct{\omega}_c &= \vct{\omega}_m + \vct{i}_{\omega} + 2K_p\,\vct{e}_{1/2},\\
        \dot{\vct{i}}_{\omega} &= 2K_i\,\vct{e}_{1/2},
    \end{align}
    i.e.\ in terms of the full cross product
    \(\vct{e}=\hat{\vct{a}}_b\times\hat{\vct{g}}_b\) the applied proportional
    correction is \(K_p\vct{e}\). When \(2K_i \le 0\) the integral accumulators
    are zeroed rather than frozen, so a temporarily zero integral gain does not
    leave a stale bias in the loop.

    The quaternion kinematics are
    \begin{equation}
        \dot{\hat{q}}
        =
        \frac12\,\hat{q}\otimes [0\;\vct{\omega}_c],
    \end{equation}
    integrated per sample with a first-order step and followed by
    renormalization using a fast inverse square root.

    \subsection{Body-to-World Rotation}

    The code rotates a body-frame vector into the Mahony world frame using the
    efficient quaternion-vector formula
    \begin{align}
        \vct{t} &= 2\,\vct{q}_v \times \vct{f}_b,\\
        \hat{\vct{f}}_w &= \vct{f}_b + q_0\vct{t} + \vct{q}_v \times \vct{t},
    \end{align}
    with \(\vct{q}_v=[q_1\;q_2\;q_3]^\top\) taken from the stored quaternion
    directly. The conjugate is exposed for callers as
    \texttt{quaternionBodyToWorld()} but is not used in the acceleration path.

    The resulting world-frame specific force is
    \begin{equation}
        \hat{\vct{f}}_w = [\hat{f}_{w,x}\;\hat{f}_{w,y}\;\hat{f}_{w,z}]^\top,
    \end{equation}
    and the vertical inertial acceleration passed downstream is
    \begin{equation}
        \hat{a}_z = \hat{f}_{w,z} - g.
    \end{equation}

    At rest and level, \(\hat{f}_{w,z}\approx g\), so \(\hat{a}_z \approx 0\).

    \subsection{Mahony Gain Adaptation}

    The wrapper can also adapt Mahony gains according to how trustworthy the
    accelerometer looks as a gravity cue. The adaptation runs at the top of
    every \texttt{updateIMU()} / \texttt{updateIMUMag()} call, on the raw body
    acceleration and against the pre-update quaternion, so the gains used by a
    sample are computed from that same sample. It computes:
    \begin{align}
        \varepsilon_{\mathrm{norm}} &=
        \frac{\left|\|\vct{a}_b\| - g\right|}{g},\\
        \varepsilon_{\mathrm{innov}} &=
        \left\|\hat{\vct{a}}_b \times \hat{\vct{g}}_b\right\|,\\
        \sigma_a &= \text{running accel sigma from the adaptive core}.
    \end{align}
    The sigma used here is the \(\sigma_a\) EMA, not the envelope amplitude of
    \cref{sec:envelope}; a non-positive or non-finite value is replaced by the
    configured reference.

    Those are normalized by configured reference scales and converted into trust
    factors
    \begin{equation}
        \rho_n = \frac{1}{1+n^2},\qquad
        \rho_e = \frac{1}{1+u^2},\qquad
        \rho_s = \frac{1}{1+s^2},
    \end{equation}
    where \(n\), \(u\), and \(s\) are the normalized norm error, innovation,
    and vertical-acceleration sigma.

    The combined accelerometer trust is
    \begin{equation}
        \gamma = \clamp(\rho_n\,\rho_e\,\rho_s,\gamma_{\min},1).
    \end{equation}

    That trust interpolates between rough-water and calm-water Mahony gains:
    \begin{align}
        2K_p(\gamma)
        &=
        2K_{p,\mathrm{rough}} +
        \gamma\bigl(2K_{p,\mathrm{calm}}-2K_{p,\mathrm{rough}}\bigr),\\
        2K_i(\gamma)
        &=
        2K_{i,\mathrm{rough}} +
        \gamma^2\bigl(2K_{i,\mathrm{calm}}-2K_{i,\mathrm{rough}}\bigr).
    \end{align}

    So the proportional term fades linearly with trust, while the integral term
    collapses even faster in rough motion. The commanded gains are then smoothed
    with the same one-pole coefficient form and written back into the Mahony
    state through \texttt{Mahony\_AHRS::init()}, which sets gains only and leaves
    the quaternion and the integral accumulators untouched. With
    \texttt{adapt\_mahony\_gains=false} the currently active gains are simply
    re-applied each sample.

% ---------------------------------------------------------------
    \section{Simulation-Backed Tuning Example}\label{sec:tuning}

    The simulation driver \texttt{tests/pii\_observer/pii\_observer-adaptive.cpp}
    defines \texttt{FusionAdapterAdaptivePIIMahony} around
    \texttt{AdaptiveVerticalPIIMahony<float,\ true,\ TrackerType::PLL>}, so the
    internal acceleration-frequency proxy there is the PLL tracker with the
    shared default tracker configuration, even though the adaptive wrapper
    itself remains generic. It sets the following values for the core observer:
    \begin{align}
        r &= 0.125,\\
        \tau_a &= 0.40~\si{s},\\
        \tau_d &= 49.0~\si{s},\\
        k_b &= 2.5\times10^{-5},\\
        \lambda_b &= 3.0\times10^{-3}.
    \end{align}

    Safety limits are
    \begin{equation}
        |a_f|\le 50,\quad |v|\le 50,\quad |p|\le 20,\quad |S|\le 200,\quad |d|\le 20,
    \end{equation}
    with
    \begin{equation}
        |b|\le 0.12.
    \end{equation}

    The adaptation layer uses
    \begin{align}
        f_{\mathrm{ref}} &= 0.12~\si{Hz},\\
        \sigma_{\mathrm{ref}} &= 1.10,\\
        \tau_{\mathrm{in}} &= 55~\si{s},\\
        \tau_{\mathrm{param}} &= 50~\si{s},\\
        c_{\min} &= 0.22,
    \end{align}
    together with the exponent set and bounds of \cref{tab:exponents}, and it is
    auto-scheduled from the accel-frequency proxy every \SI{0.5}{s}. The
    fallback confidences are \(0.52\) with a coarse estimate, \(0.62\) from the
    coarse-schedule floor, and \(0.82\) when locked, with a coarse blend of
    \(0.48\).

    \begin{table}[t]
        \centering
        \caption{Scheduling exponents and bounds used in the simulation}
        \label{tab:exponents}
        \small
        \begin{tabular}{@{}lrrl@{}}
            \toprule
            Parameter & \(\rho_f\) exp. & \(\rho_\sigma\) exp. & Bounds \\
            \midrule
            \(r\)       & \(0.23\)  & \(0.65\)  & \([0.06,\,0.26]\) \\
            \(\tau_a\)  & \(-0.25\) & \(-1.20\) & \([0.15,\,0.90]~\si{s}\) \\
            \(\tau_d\)  & \(-0.03\) & \(-0.01\) & \([44,\,58]~\si{s}\) \\
            \(k_b\)     & \(0.02\)  & \(0.08\)  & \([5\times10^{-6},\,6\times10^{-5}]\) \\
            \bottomrule
        \end{tabular}
    \end{table}

    The envelope estimator is enabled with the OU-aligned values of
    \cref{sec:envelope}: \(\sigma_n=\SI{0.12}{m/s^2}\), \(K_{\mathrm{per}}=2.0\),
    \(\tau_f=\SI{1.0}{s}\), \(c_\tau=1.38\), \(c_\sigma=0.90\), an adaptation
    time constant of \SI{1.8}{s}, \(\tau\in[0.02,3.00]~\si{s}\), and
    \(\sigma_{\max}=6.0\).

    The simulation-backed Mahony tuning uses
    \begin{align}
    (2K_p)_{\mathrm{base}} &= 1.70,\\
    (2K_i)_{\mathrm{base}} &= 0.090,\\
    (2K_p)_{\mathrm{calm}} &= 1.50,\\
    (2K_p)_{\mathrm{rough}} &= 1.20,\\
    (2K_i)_{\mathrm{calm}} &= 0.100,\\
    (2K_i)_{\mathrm{rough}} &= 0.070.
    \end{align}

    The trust references are
    \begin{align}
        \sigma_{\mathrm{ref}}^{\mathrm{Mahony}} &= 0.45~\si{m/s^2},\\
        \varepsilon_{\mathrm{norm,ref}} &= 0.12,\\
        \varepsilon_{\mathrm{innov,ref}} &= 0.18,
    \end{align}
    with Mahony gain-schedule smoothing time constant
    \begin{equation}
        \tau_g = 1.0~\si{s},
    \end{equation}
    and minimum accelerometer trust
    \begin{equation}
        \gamma_{\min} = 0.65.
    \end{equation}

    The high trust floor is what keeps this tuning band narrow: with
    \(\gamma\in[0.65,1]\) the scheduled \(2K_p\) only ever moves between
    \(1.40\) and \(1.50\), so the Mahony gain schedule trims the attitude loop
    rather than reshaping it.

    \subsection{Frame Mapping in the Simulation Adapter}

    The simulation runner provides body NED data, and the adapter converts it to
    the Mahony body convention. Inertial and magnetic channels use
    \emph{different} mappings, and the code says so explicitly, because sharing
    one of them causes a large yaw error. For gyro and accelerometer,
    \begin{equation}
    (x,y,z)_{\mathrm{NED}}
    \mapsto
    (-y,-x,-z)_{\mathrm{Mahony}},
    \end{equation}
    which is the \(Z\)-up conversion \((x,y,z)\mapsto(y,x,-z)\) followed by
    negation of the two horizontal axes. At rest and level the runner's
    accelerometer reads \((0,0,-g)_{\mathrm{NED}}\), which maps to
    \((0,0,+g)_{\mathrm{Mahony}}\) as the Mahony convention requires.

    For the magnetometer,
    \begin{equation}
    (x,y,z)_{\mathrm{NED}}
    \mapsto
    (x,-y,-z)_{\mathrm{Mahony}},
    \end{equation}
    so magnetic north stays on \(+X\), which is what this Mahony mag path
    expects. The latest magnetometer sample is cached and reused at every IMU
    step until a new one arrives, since the mag ODR (\SI{20}{Hz} in the runs
    below) is well under the \SI{200}{Hz} IMU rate.

    For RMS reporting back in the simulator convention, the adapter decodes the
    Mahony quaternion to nautical Euler angles. Yaw needs one further
    correction: the reference trace is reported in the magnetic frame as
    \(\psi_{\mathrm{ref}}+\delta\) while the decoded Mahony yaw lands on
    \(\psi_{\mathrm{ref}}-\delta\), so the adapter reports
    \begin{equation}
        \psi_{\mathrm{sim}} = \mathrm{wrap}\bigl(\psi_{\mathrm{Mahony}}
        + 2\delta\bigr),
    \end{equation}
    with \(\delta\) the simulated declination. Without magnetometer aiding the
    reported yaw is forced to zero rather than scored.

    The adapter also republishes observer telemetry for the plots: the active
    \(r\) and \(k_b\) as the tuning traces, and \emph{envelope}
    \(\tau,\sigma,f\), and variance as the envelope traces. The older
    \(\sigma_{a}\) filter states, \(f_{\mathrm{disp}}\) filter state, core
    \texttt{accel\_var}, and \(\tau_a\) are deliberately no longer used for the
    wave-height envelope.

    \subsection{Quality Gates Used in the Example}

    The simulation code checks the last \SI{900}{s} of each \SI{1200}{s} run and
    applies quality gates:
    \begin{itemize}[leftmargin=1.4em]
        \item Z RMS \(< 9.0\%\) of \(H_s\) for JONSWAP cases,
        \item Z RMS \(< 9.5\%\) of \(H_s\) for PM-Stokes cases,
        \item yaw RMS \(< 10.9^\circ\) when magnetometer is used.
    \end{itemize}
    A record shorter than the scoring window is reported as skipped rather than
    scored, so the silence of a short record cannot read as a pass.

    These are not part of the observer itself, but they are useful context for the
    intended tuning regime in the supplied simulation.

    The window is the one the OU-II/OU-III simulators gate on and the one the
    four-method comparison of the OU-III manuscript scores, so this observer is
    described over the same stretch of the replay as the estimators it is
    compared against. Each threshold is a regression sentinel for the
    deterministic realization -- the worst value this observer currently produces
    across the eight scored records plus about half a percent -- rather than an
    acceptance criterion. The worst record for the Z gates is the roughest sea
    (\(H_s=\SI{8.5}{m}\)) under both spectra, which is where the sentinels moved
    once scheduling the PII pole rate off accel sigma removed the calm-sea error.

    \begin{table}[t]
        \centering
        \caption{Representative tuning from the simulation adapter}
        \label{tab:tuning}
        \small
        \begin{tabular}{@{}lll@{}}
            \toprule
            Parameter & Value & Role \\
            \midrule
            \(r\) & 0.125 & base repeated pole rate \\
            \(\tau_a\) & \SI{0.40}{s} & accel LPF time constant \\
            \(\tau_d\) & \SI{49}{s} & slow trend time constant \\
            \(k_b\) & \(2.5\times10^{-5}\) & bias learning gain \\
            \(\lambda_b\) & \(3.0\times10^{-3}\) & bias leak \\
            \(f_{\mathrm{ref}}\) & \SI{0.12}{Hz} & reference scheduling frequency \\
            \(\sigma_{\mathrm{ref}}\) & 1.10 & reference accel sigma \\
            \(\tau_{\mathrm{in}}\) & \SI{55}{s} & scheduling input smoothing \\
            \(\tau_{\mathrm{param}}\) & \SI{50}{s} & scheduled parameter smoothing \\
            \((2K_p)_{\mathrm{base}}\) & 1.70 & Mahony base proportional gain \\
            \((2K_i)_{\mathrm{base}}\) & 0.090 & Mahony base integral gain \\
            \(\gamma_{\min}\) & 0.65 & minimum accelerometer trust \\
            \(g\) & \SI{9.80665}{m/s^2} & gravity magnitude \\
            \bottomrule
        \end{tabular}
    \end{table}

% ---------------------------------------------------------------
    \section{Overall Interpretation}

    The complete method is a cascaded observer for vertical motion in rough seas:
    \begin{enumerate}[leftmargin=1.5em]
        \item Mahony AHRS estimates orientation from gyroscope integration corrected by
        gravity and optional magnetometer aiding.
        \item Body-frame specific force is rotated into a \(Z\)-up world frame and
        converted to vertical inertial acceleration by subtracting gravity.
        \item A policy-selected frequency tracker provides a dominant motion-frequency
        proxy and confidence information, blended toward a coarse estimate when
        the tracker is not locked.
        \item An OU-style envelope estimator turns that frequency and a
        period-scaled acceleration variance into a sea-state amplitude and
        correlation time, which are the preferred scheduling inputs and the
        source of the reported wave-height and vertical-speed scales.
        \item A vertical PII observer integrates the vertical acceleration to
        displacement while using repeated-pole damping, displacement-integral
        feedback, and an optional very-slow bias channel to prevent unbounded
        drift.
        \item The observer bandwidth and bias-learning rates are scheduled from the
        estimated motion frequency and sea-state amplitude, so calmer conditions
        permit tighter tracking while rougher conditions slow the loop and reduce
        sensitivity to spurious bias inference.
        \item The Mahony gains are also reduced when accelerometer trust is poor, so
        tilt correction does not chase wave-induced specific force as if it were
        gravity.
    \end{enumerate}

    This design avoids a full stochastic state-space model yet still introduces
    physically meaningful structure: attitude resolves gravity, the frequency
    proxy provides a representative time scale, and the adaptive PII core shapes
    heave dynamics to that estimated motion regime.

    \section{Results Plots}

    \begin{figure}[!t]\centering
    \IncludePGF{w3d_nonkalman_pmstokes_medium.pgf}
    \caption{Adaptive PII/Mahony attitude (roll, pitch, yaw) --- PM Stokes, medium $H_s$.}
    \label{fig:w3d_nonkalman_pmstokes_medium_angles}
    \end{figure}
    \begin{figure}[!t]\centering
    \IncludePGF{w3d_nonkalman_pmstokes_medium_zkin.pgf}
    \caption{Adaptive PII/Mahony vertical kinematics (displacement, velocity, acceleration with learned envelopes) --- PM Stokes, medium $H_s$.}
    \label{fig:w3d_nonkalman_pmstokes_medium_zkin}
    \end{figure}
    \begin{figure}[!t]\centering
    \IncludePGF{w3d_nonkalman_pmstokes_medium_tuner.pgf}
    \caption{Adaptive PII/Mahony frequency and auto-scheduling traces (envelope $f$, $\sigma$, $\tau$, and active PII $r$) --- PM Stokes, medium $H_s$.}
    \label{fig:w3d_nonkalman_pmstokes_medium_tuner}
    \end{figure}

% ---------------------------------------------------------------
    \section{Summary}

    The code implements a compact, layered vertical-motion estimation stack.

    First, \texttt{VerticalPIIObserver} is a repeated-real-pole PII observer with
    optional very-slow bias estimation. It is intentionally conservative,
    non-oscillatory, and inexpensive to run.

    Second, \texttt{AdaptiveVerticalPII} adds acceleration sigma estimation, an
    OU-style sea-state envelope estimator, a policy-selected frequency tracker,
    and runtime scheduling of the observer parameters. External scheduling
    inputs are honored when supplied, the envelope is the preferred internal
    source, and the plain sigma EMA is the fallback until the envelope is ready.

    Third, \texttt{AdaptiveVerticalPIIMahony} uses Mahony AHRS to supply the
    required vertical inertial acceleration and can also adapt Mahony gains from
    accelerometer trust and sea-state roughness indicators.

    Together, these layers form a practical embedded heave-estimation pipeline
    that is easy to tune, robust to moderate motion variability, and compatible
    with several lightweight frequency-tracker backends.

\end{document}
